% =====================================================================
%  Project 3 - Matrix multiplication, inversion, and photo filters
%  Overleaf: New Project > Upload Project > chọn file bao_cao.zip
%  Compiler: pdfLaTeX (mặc định của Overleaf). Biên dịch: Recompile.
%  Mỗi người chỉ sửa file của mình trong sections/ (tên người viết ghi ở đầu file).
% =====================================================================
\documentclass[11pt,a4paper]{article}
\usepackage[utf8]{inputenc}
% Tiếng Việt cho họ tên: Overleaf có sẵn gói T5. Máy không có thì dùng T1.
\IfFileExists{t5enc.def}{\usepackage[T5,T1]{fontenc}}{\usepackage[T1]{fontenc}}
\IfFileExists{lmodern.sty}{\usepackage{lmodern}}{}
\usepackage[margin=2.3cm]{geometry}
\usepackage{amsmath,amssymb,amsthm}
\usepackage{graphicx,float,booktabs,array,caption,subcaption}
\usepackage{xcolor,listings}
\usepackage[hidelinks]{hyperref}
\usepackage{cleveref}

\graphicspath{{figures/}}
\newtheorem{definition}{Definition}
\newtheorem{theorem}{Theorem}

% ---- ghi chú việc còn thiếu (xóa hết trước khi nộp) ----
\newcommand{\todo}[1]{\textcolor{red}{\textbf{[TODO: #1]}}}
\newcommand{\writer}[1]{\marginpar{\footnotesize\textcolor{gray}{#1}}}
% Trước khi nộp: bỏ dấu % ở dòng dưới để ẩn mọi TODO và tên người viết
% \renewcommand{\todo}[1]{}\renewcommand{\writer}[1]{}

% ---- định dạng code Python ----
\lstset{language=Python, basicstyle=\ttfamily\footnotesize, numbers=left,
  numberstyle=\tiny\color{gray}, keywordstyle=\color{blue!70!black}\bfseries,
  commentstyle=\color{green!40!black}\itshape, stringstyle=\color{red!60!black},
  frame=single, breaklines=true, columns=fullflexible, showstringspaces=false,
  upquote=true, inputencoding=utf8, extendedchars=true,
  literate={≥}{{$\geq$}}1 {≤}{{$\leq$}}1 {→}{{$\to$}}1 {−}{{-}}1}

\begin{document}

% ============================ TRANG BÌA ============================
\begin{titlepage}
\centering
{\large HO CHI MINH CITY UNIVERSITY OF TECHNOLOGY\par}
{\large FACULTY OF APPLIED SCIENCE\par}   % \todo{kiểm tra tên khoa}
\vspace{2.5cm}
{\Large LINEAR ALGEBRA (MT1007)\par}
\vspace{0.8cm}
{\LARGE\bfseries Project 3\\[4pt] Matrix Multiplication, Inversion,\\ and Photo Filters\par}
\vspace{2cm}
\begin{tabular}{lll}
\toprule
 & Full name & Student ID \\
\midrule
1 & Ho ten TV1 (leader) & 0000000 \\
2 & Ho ten TV2 & 0000000 \\
3 & Ho ten TV3 & 0000000 \\
4 & Ho ten TV4 & 0000000 \\
5 & Ho ten TV5 & 0000000 \\
6 & Ho ten TV6 & 0000000 \\
7 & Ho ten TV7 & 0000000 \\
\bottomrule
\end{tabular}

\vspace{1.2cm}
Class: \todo{lớp} \qquad Instructor: \todo{tên giảng viên}
\vfill
{Ho Chi Minh City, \today\par}
\end{titlepage}

\tableofcontents
\newpage

% =====================================================================
%  1. PROBLEM STATEMENT  -  Người viết: TV5
%  Đề yêu cầu: vấn đề đặt ra, tầm quan trọng, ứng dụng.
%  Tự viết bằng lời của mình. Mọi ý lấy từ nguồn ngoài phải \cite{...}.
%  Độ dài gợi ý: 1 - 1,5 trang.
% =====================================================================
\section{Problem statement}\label{sec:problem}\writer{TV5}

\subsection{The problem}
% Gợi ý: ảnh số là gì (lưới pixel, mỗi pixel 3 số R, G, B trong 0..255);
% bài toán: biến đổi màu cả ảnh bằng cách nhân ma trận 3x3 với vector màu
% của từng pixel; câu hỏi chính: phép lọc nào đảo ngược được, vì sao.
\todo{mô tả bài toán}

\subsection{Why it matters}
% Gợi ý: một ví dụ trực quan về phép nhân ma trận, định thức, nghịch đảo,
% số điều kiện; liên hệ giữa "khả nghịch về lý thuyết" và "khôi phục được thật".
\todo{tầm quan trọng}

\subsection{Applications}
% Gợi ý: bộ lọc ảnh trên điện thoại / mạng xã hội, chuyển không gian màu,
% chỉnh cân bằng trắng, hỗ trợ người mù màu, nén ảnh... (có trích dẫn).
\todo{ứng dụng}

\subsection{Objectives and structure of the report}
\todo{mục tiêu nhóm đặt ra và tóm tắt các phần sau}

% =====================================================================
%  2. THEORETICAL FOUNDATION
%  Đề yêu cầu: kiến thức cần để giải bài toán, PHẢI tự nghiên cứu và tự viết,
%  KHÔNG sao chép (kể cả từ file lời giải của nhóm, sách, web hay AI).
%  Mỗi người viết một mục, mỗi mục có ít nhất 1 ví dụ tính tay nhỏ.
% =====================================================================
\section{Theoretical foundation}\label{sec:theory}

% ---------------------------------------------------------------------
\subsection{Color images as arrays and linear color filters}\label{sec:theory-filters}
\writer{TV2}   % Hạn nháp: N6
% Gợi ý nội dung:
%  - ảnh = mảng m x n x 3; kiểu uint8, vì sao cần đổi sang số thực (Q1)
%  - pixel p = [R;G;B]; bộ lọc tuyến tính p' = M p; vì sao (M p^T)^T = p M^T
%    cho phép tính cả ảnh bằng một phép nhân ma trận
%  - ý nghĩa từng hàng / cột của M; ma trận đường chéo (Saturate), ma trận
%    hoán vị (Permute), phép chiếu (Red, DeleteGreen, Gray)
%  - phép biến đổi KHÔNG tuyến tính: đảo màu (affine), gamma (Task 9, 14)
%  - làm tròn và cắt về [0, 255] sau khi lọc
\todo{viết mục 2.1}

% ---------------------------------------------------------------------
\subsection{Determinant, inverse matrix and invertible filters}\label{sec:theory-inverse}
\writer{TV7}   % Hạn nháp: N6
% Gợi ý nội dung:
%  - định nghĩa ma trận nghịch đảo; điều kiện khả nghịch (det khác 0,
%    các hàng/cột độc lập tuyến tính)
%  - cách tính det và nghịch đảo ma trận 3x3 (ví dụ tay với UserMatrix)
%  - ý nghĩa: bộ lọc khả nghịch thì về lý thuyết đảo ngược được (Q5)
%  - ma trận suy biến (Gray, Red, DeleteGreen) làm mất thông tin gì
\todo{viết mục 2.2}

% ---------------------------------------------------------------------
\subsection{Eigenvalues, conditioning and numerical inversion}\label{sec:theory-cond}
\writer{TV3}   % Hạn nháp: N6
% Gợi ý nội dung:
%  - trị riêng, vector riêng; phân tích UserMatrix (vì sao nó kéo màu về xám)
%  - chuẩn ma trận, số điều kiện cond(M); sai số đầu vào bị khuếch đại
%    tối đa cond(M) lần
%  - vì sao det nhỏ chưa đủ để kết luận, cond mới là thước đo đúng
%  - áp vào SepiaMatrix: các hàng gần tỉ lệ -> gần suy biến (Q6)
\todo{viết mục 2.3}

% =====================================================================
%  3. CODE  -  Người viết: TV1
%  Bản nháp: TV1 đọc lại, sửa bằng lời của mình trước khi gửi TV5.
% =====================================================================
\section{Code}\label{sec:code}\writer{TV1}

The project was implemented in Python instead of MATLAB, as recommended in the
assignment, because Python is open source and free of licensing issues. The
whole program is a single script, \texttt{lab03.py}, listed in
Appendix~\ref{app:code}. It keeps the variable names of the original MATLAB
project (\texttt{ImJPG}, \texttt{redChannel}, \texttt{ImJPG\_Gray}, \dots) so
that each result can be matched with the corresponding task.

\subsection{Environment}
The code needs Python~3 and three libraries:
\begin{itemize}
  \item \textbf{NumPy}~\cite{harris2020numpy}: arrays and matrix operations
        (\texttt{@}, \texttt{linalg.det}, \texttt{linalg.inv}, \texttt{linalg.cond});
  \item \textbf{Pillow}: reading the image file \texttt{daylilies.jpg};
  \item \textbf{Matplotlib}~\cite{hunter2007matplotlib}: displaying the images.
\end{itemize}
They are installed and the program is run with:
\begin{lstlisting}[language=bash,numbers=none]
pip install numpy matplotlib pillow
python lab03.py      # daylilies.jpg must be in the same folder
\end{lstlisting}

\subsection{Program structure}
The script has two parts. The first part imports the libraries and defines
three shared functions used by every task (\texttt{to\_uint8},
\texttt{apply\_filter\_loop}, \texttt{show}). The second part contains the
14 tasks in order; each task starts with a cell marker \texttt{\# \%\% Task k},
so it can be run on its own in VS Code or Spyder. The answers to questions
Q1--Q6 are written as comments directly under the related task. All figures
are shown at once by \texttt{plt.show()} at the end of the file.

Since the seven members wrote their tasks in separate files
(\texttt{part1.py}, \dots, \texttt{part7.py}), a small helper script,
\texttt{ghep\_file.py}, merges them into \texttt{lab03.py}: it keeps the shared
part only once, puts the tasks in the order 1 to 14 (Task 12 needs a variable
from Task 11, and Task 13 one from Task 5), and checks that every variable
required by the assignment exists.

Table~\ref{tab:matlab-python} lists the main MATLAB commands of the project
and the Python code used instead.

\begin{table}[H]\centering\small
\caption{MATLAB commands of the project and their Python equivalents.}
\label{tab:matlab-python}
\begin{tabular}{lll}
\toprule
MATLAB & Python & Remark \\
\midrule
\texttt{imread('daylilies.jpg')} & \texttt{np.array(Image.open(...).convert("RGB"))} & array $m\times n\times 3$, uint8 \\
\texttt{[m,n,l] = size(ImJPG)} & \texttt{m, n, l = ImJPG.shape} & \\
\texttt{ImJPG(:,:,1)} & \texttt{ImJPG[:, :, 0]} & indices start at 0 \\
\texttt{A*B} & \texttt{A @ B} & \texttt{A * B} is element-wise \\
\texttt{for i=1:m ... end} & \texttt{for i in range(m):} & loop over pixels \\
\texttt{reshape(x,3,1)} & \texttt{x.reshape(3, 1)} & column vector \\
\texttt{double(x)} & \texttt{x.astype(np.float64)} & \\
\texttt{uint8(x)} & \texttt{to\_uint8(x)} & own function \\
\texttt{inv, det, cond} & \texttt{np.linalg.inv, det, cond} & \\
\texttt{figure; imshow; title} & \texttt{show(img, title)} & own function \\
\bottomrule
\end{tabular}
\end{table}

\subsection{Key functions}

\paragraph{\texttt{to\_uint8(x)}.}
MATLAB's \texttt{uint8} rounds a number to the nearest integer and saturates
it to $[0,255]$, e.g.\ \texttt{uint8(260)} gives 255. In NumPy,
\texttt{.astype(np.uint8)} does not saturate but wraps around modulo 256
(260 becomes 4), which would create wrong colored spots in bright areas.
The function therefore rounds halves away from zero, as MATLAB does, and
clips the result explicitly:
\begin{lstlisting}
def to_uint8(x):
    x = np.sign(x) * np.floor(np.abs(x) + 0.5)
    return np.clip(x, 0, 255).astype(np.uint8)
\end{lstlisting}

\paragraph{\texttt{apply\_filter\_loop(img, M)}.}
A filter multiplies the color vector $p=[R;G;B]$ of every pixel by a
$3\times3$ matrix $M$. Following the project text, this is done with two
nested loops over the pixels: for each pixel the color vector is taken,
converted to floating point (MATLAB \texttt{double}), reshaped to a
$3\times1$ column (MATLAB \texttt{reshape}), multiplied by the filter matrix
and converted back with \texttt{to\_uint8} (MATLAB \texttt{uint8}):
\begin{lstlisting}
def apply_filter_loop(img, M):
    m, n, l = img.shape
    out = np.zeros((m, n, l), dtype=np.uint8)
    for i in range(m):
        for j in range(n):
            PixelColor = img[i, j, :].astype(np.float64).reshape(3, 1)
            out[i, j, :] = to_uint8(M @ PixelColor).ravel()
    return out
\end{lstlisting}
The pixel values are converted to floating point before the product, so that
no intermediate result is rounded or overflows. All matrix filters in
\texttt{lab03.py} (Tasks 4--8 and 10--13) use this single function. Filtering
the $308\times322$ image takes about one second.

\paragraph{\texttt{show(img, title)}.}
Opens a new figure and displays the image with a title. A single-channel image
(for example \texttt{redChannel}) is shown in gray levels.

\subsection{Differences from the MATLAB version and verification}
Besides the indexing from 0 and the operator \texttt{@}, two tasks needed extra
care. In Task~9 the negative is computed as \texttt{255 - ImJPG}, which never
leaves $[0,255]$. In Task~14 MATLAB adds 30 to a \texttt{uint8} image with
saturation, while NumPy would wrap around (250 + 30 = 24); the addition is
therefore done in floating point and followed by \texttt{to\_uint8}.

The results were checked against the values expected from the project:
the image size is $m=308$, $n=322$, $l=3$; $\det(\texttt{UserMatrix}) = 0.3025$;
restoring \texttt{ImJPG\_User} with the inverse filter differs from the
original image by at most 1 (out of 255); $\det(\texttt{SepiaMatrix})\approx 1.21\times10^{-7}$
and $\operatorname{cond}(\texttt{SepiaMatrix})\approx 16\,659$. The figures of
Tasks 3 and 7 match Fig.~5 and Fig.~6 of the project text.

% =====================================================================
%  4. ILLUSTRATIVE EXAMPLE
%  Đề: dùng ví dụ trong tài liệu để kiểm tra code đúng; XÂY DỰNG hoặc TÌM
%  thêm dữ liệu làm ví dụ.
%  Hình và bảng do vi_du_them.py tạo trong figures/ - chạy lại script là cập nhật.
% =====================================================================
\section{Illustrative examples}\label{sec:examples}

% ---------------------------------------------------------------------
\subsection{The example in the material: \texttt{daylilies.jpg}}\label{sec:ex-material}
\writer{TV4}   % Hạn nháp: N6
% Gợi ý: mỗi task 1-2 câu + câu trả lời Q1-Q6 (viết lại, không dán từ code),
% so khớp với Fig. 5, Fig. 6 của đề, số liệu m, n, l, det, cond.
\todo{trình bày kết quả 14 task}

\IfFileExists{figures/luoi_daylilies.png}{
\begin{figure}[H]\centering
\includegraphics[width=\textwidth]{luoi_daylilies.png}
\caption{Filters of Tasks 4--14 applied to the image in the material.}\label{fig:grid-daylilies}
\end{figure}
\begin{figure}[H]\centering
\includegraphics[width=\textwidth]{khoi_phuc_daylilies.png}
\caption{Inverting the User filter (Task 12) and the sepia filter (Task 13).}\label{fig:restore-daylilies}
\end{figure}}{\todo{chạy vi\_du\_them.py khi đã có 00\_DeBai/daylilies.jpg}}

\begin{table}[H]\centering
\caption{Determinant and condition number of every filter matrix.}\label{tab:matrices}
\IfFileExists{figures/bang_ma_tran.tex}{\input{figures/bang_ma_tran.tex}}{\todo{chạy vi\_du\_them.py}}
\end{table}

% ---------------------------------------------------------------------
\subsection{Additional data built and collected by the group}\label{sec:ex-extra}
\writer{TV6}   % Hạn nháp: N6
% Gợi ý:
%  - Ảnh tổng hợp (bang_mau.png): dải màu thuần + dải xám do nhóm tự xây,
%    cho thấy rõ từng ma trận làm gì với từng màu, vì sao UserMatrix giữ nguyên xám.
%  - 2-3 ảnh nhóm tự chụp / tìm (ghi nguồn + giấy phép), bỏ vào 05_DuLieuThem/.
%  - Nhận xét bảng khôi phục: User khôi phục gần đúng mọi ảnh, Sepia thì không;
%    ảnh nào bị cắt giá trị nhiều nhất, vì sao.
\todo{mô tả dữ liệu thêm và nhận xét}

\IfFileExists{figures/luoi_bang_mau.png}{
\begin{figure}[H]\centering
\includegraphics[width=0.55\textwidth]{bang_mau.png}
\caption{Synthetic color chart built by the group: pure hues from dark to bright, and a gray ramp.}\label{fig:chart}
\end{figure}}{}

\IfFileExists{figures/hinh_du_lieu_them.tex}{\input{figures/hinh_du_lieu_them.tex}}{}

\begin{table}[H]\centering\small
\caption{Recovering images with the inverse filters (errors on the 0--255 scale) and share of channel values clipped to $[0,255]$.}\label{tab:restore}
\IfFileExists{figures/bang_khoi_phuc.tex}{\input{figures/bang_khoi_phuc.tex}}{\todo{chạy vi\_du\_them.py}}
\end{table}

% =====================================================================
%  5. COMMENTS ON THE TOPIC  -  Người viết: TV7
%  Đề: đã làm được gì, chưa làm được gì, hướng phát triển.
%  Thành thật: nêu cả hạn chế (vd. chỉ xét bộ lọc 3x3 trên từng pixel).
% =====================================================================
\section{Comments on the topic}\label{sec:comments}\writer{TV7}

\subsection{What has been accomplished}
\todo{liệt kê kết quả đạt được, đối chiếu với yêu cầu đề}

\subsection{Limitations}
\todo{những gì chưa làm được / chưa chính xác}

\subsection{Possible further development}
% Ý tưởng: bộ lọc có độ lệch (affine 4x4), khôi phục bằng bình phương tối
% thiểu hoặc chính quy hóa, bộ lọc theo vùng lân cận (tích chập: làm mờ, làm nét),
% không gian màu khác (YCbCr, HSV), giao diện chọn bộ lọc...
\todo{hướng phát triển}


\clearpage
\appendix
\section{Full source code \texttt{lab03.py}}\label{app:code}
\IfFileExists{code/lab03.py}
  {\lstinputlisting{code/lab03.py}}
  {\todo{TV1 chép lab03.py (bản cuối) vào thư mục code/}}

\section{Work distribution}
% TV1 cập nhật bảng này theo phân công thật
\begin{tabular}{>{\raggedright}p{3.2cm}p{5.6cm}p{5.6cm}}
\toprule
Member & Code & Report \\
\midrule
TV1 (leader) & Task 1--2, merging, final version, submission & Sec.~\ref{sec:code}, PDF export \\
TV2 & Task 4, 11; shared functions; Q1, Q4 & Sec.~\ref{sec:theory-filters} \\
TV3 & Task 5, 13; Q6 & Sec.~\ref{sec:theory-cond} \\
TV4 & Task 3, 6; Q2; first test run & Sec.~\ref{sec:ex-material} \\
TV5 & Task 7, 14; progress tracking & Sec.~\ref{sec:problem}, references \\
TV6 & Task 9, 10; Q3; final test run & Sec.~\ref{sec:ex-extra} \\
TV7 & Task 8, 12; Q5 & Sec.~\ref{sec:theory-inverse}, \ref{sec:comments} \\
\bottomrule
\end{tabular}

\clearpage
% ======================= 6. REFERENCES (TV5) =======================
\nocite{zemlyanova2020,strang2016}  % tài liệu chính của đề, luôn liệt kê
\bibliographystyle{plain}
\bibliography{refs}

\end{document}
